%======================================================================
\section{자료와 방법}

\begin{yoji}
아비트럼 원의 공개 로그에서 허락과 송금을 읽어, 동결 직전 12개월의 사건으로 동결마다 그 코호트(같은 기간 평가와 W1은 계약 6,483개, W0은 6,343개)를 둘러싼 그래프를 만들었어요. 점수를 동결한 뒤 석 달 동안 생긴 새 허락 쌍과 새 송금자(스펜더 라벨), 그리고 GMX V2 계약 트레이더의 청산(트레이더 라벨)으로 채점해요. 무엇을 언제 정했는지도 순서대로 밝혀요.
\end{yoji}

\subsection{자료와 코호트}

모든 사건은 이더리움 위의 낙관적 롤업인 아비트럼 원에서 와요~\cite{kalodner2018arbitrum}. 공개 BigQuery 표 \texttt{goog\_blockchain\_arbitrum\_one\_us.logs}(프로젝트 \texttt{bigquery-public-data})에서 읽었어요. 추출 범위는 2023년 10월 1일부터 2026년 6월 30일까지이고, 두 번째 창의 라벨을 위해 2026년 7월부터 9월까지도 읽었어요.

이 원고의 모든 점수는 동결 직전 12개월, 곧 그 점수의 \textbf{점수 기간}의 사건으로 만들어요.
\begin{itemize}
\item 같은 기간 평가와 두 번째 창: 2025년 7월 1일부터 2026년 6월 30일까지. 끝 시각인 23:59:59~UTC를 $\Tobs$로 써요.
\item 첫 번째 창: 2025년 4월 1일부터 2026년 3월 31일까지.
\end{itemize}
1년짜리 기간은 각 계약을 동결 전 1년의 활동으로 묘사하고, 두 동결에 같은 길이의 기간을 줘요.

읽는 이벤트는 세 가지예요. ERC-20 \texttt{Approval}과 \texttt{Transfer} 로그는 토픽이 정확히 셋인 것만 읽어요. 같은 서명을 쓰는 ERC-721 로그를 빼기 위해서예요. 그리고 아비트럼 원의 무기한 선물 거래소 GMX~V2의 \texttt{PositionDecrease} 이벤트를 읽어요~\cite{gmxnddocs}. 토픽, 풀어낸 필드, 라벨 식, 질의는 보조 자료 S1 Text에 있어요.

\paragraph{계약.}
고정한 블록에서 \texttt{eth\_getCode}가 비어 있지 않은 코드를 돌려주면 그 주소를 계약으로 봐요. 코드가 비어 있거나, 코드가 EIP-7702 위임 표시뿐인 주소~\cite{buterin2024eip7702}는 외부 소유 계정(EOA)이에요. 분석 계획의 후보 스펜더는 $\Tobs$까지 추출한 로그에서 서로 다른 주인 셋 이상에게 0이 아닌 허락을 받은 모든 주소이고, 29,022개였어요. 블록 513,216,506(2026년 10월 9일)에서 그 가운데 27,844개는 계약 코드를 가졌고, 1,121개는 코드가 없었고, 57개는 위임 표시를 가졌어요. 점수를 매기고 평가하는 대상은 계약뿐이에요. 개인의 것일 수 있는 EOA는 주인, 보낸 쪽, 받는 쪽으로만 나오고, 순위를 매기거나 보고하지 않아요.

\paragraph{토큰과 가격.}
서로 다른 토큰의 금액은 한 단위로 값을 매긴 뒤에야 더할 수 있어요. 그래서 다섯 토큰(WETH, USDC, USDT0, ARB, WBTC)의 허락과 송금만 남기고, 모든 금액을 DefiLlama의 하루 종가로 달러 값을 매겨요~\cite{defillamanddocs}. 토큰은 두 조건을 모두 만족할 때 받아들였어요.
\begin{itemize}
\item Chainlink가 아비트럼 원에 내놓은 데이터 피드 목록에, 그 토큰의 시장 가격 피드가 있고 그 시장 위험 등급이 낮음이나 중간이어야 해요~\cite{chainlinknddatafeeds}. 그 토큰이 1:1로 감싸는 자산의 피드나, 그 토큰이 정식 브리지로 옮겨 온 이더리움 원본 토큰의 피드여도 돼요.
\item 그 토큰이 자료에 나와야 해요.
\end{itemize}
목록이 아비트럼 원 주소에 연결한 토큰 133개 가운데 64개가 자격을 얻었고, 그 가운데 46개는 하루 가격이 있어요. 그중 $\Tobs$의 점수 기간에 행이 가장 많은 다섯을 남겼어요. 이 다섯이 그 기간 모든 토큰 행의 79.40\%예요. 점수 기간과 라벨 기간의 모든 날에 다섯 토큰 모두 가격이 있어요. 토큰마다 이 날들 가운데 20일을 골라 체인 위에서 읽은 Chainlink 값과 견줬더니, DefiLlama 가격은 토큰마다 중앙값으로 0.010\%에서 0.056\%만큼 달랐고, 가장 큰 차이는 0.24\%였어요.

송금은 그날의 가격으로, 마지막 허락은 기준 시점이 되는 날의 가격으로 값을 매겨요. 허락 화살표는 기준 시점에 살아 있는 권한 부여를 나타내기 때문이에요. 다섯 토큰만 남기면, 허락하거나 돈을 낸 적 없는 주인의 이름으로 스팸 토큰과 주소 중독(address poisoning) 토큰이 남길 수 있는 로그도 함께 빠져요. 이더 자체는 \texttt{Transfer} 로그 없이 움직이므로 그래프에는 WETH로만 들어가요.

\paragraph{코호트.}
\begin{itemize}
\item \textbf{매칭 코호트:} $\Tobs$의 점수 기간 안에서 서로 다른 주인 셋 이상이 다섯 토큰 가운데 하나로 0이 아닌 금액을 허락한 계약 6,483개예요. 계획의 계약 27,844개 가운데, 같은 규칙을 다섯 토큰과 그 기간에 적용해도 남는 계약이에요.
\item \textbf{W0 코호트:} 첫 번째 동결은 자기 코호트를 따로 가져요. 동결 뒤의 어떤 사건도 코호트 선택에 끼어들지 않게, 같은 규칙으로 자기 점수 기간에서 뽑았어요. 6,343개이고, 그 가운데 5,383개는 매칭 코호트에도 있어요.
\item \textbf{에고 사건:} 추출은 계획의 코호트 계약이 한쪽 끝에 있는 사건만 남겼어요. 그 계약이 주인이거나 스펜더인 허락, 그 계약이 보낸 쪽이거나 받는 쪽인 0이 아닌 송금이에요. 동결마다 자기 코호트의 사건을 자기 점수 기간과 라벨 기간에서 써요. 그래서 그래프는 코호트를 가운데 둔 한 걸음 이웃(에고 네트워크)이고, 화살표 두 개 이상의 경로는 허락을 주기도 하고 받기도 하는 코호트 계약을 거칠 때만 생겨요. $\Tobs$의 점수 기간에 있는 계획 코호트의 에고 사건 가운데, 분석은 허락 85,986,562건과 송금 442,722,039건을 남겨요. 각각 63.3\%와 76.2\%예요.
\item \textbf{스펜더 코호트:} 동결마다, 점수 기간 안에 다섯 토큰 가운데 하나로 허락을 받았고 그 허락이 동결 시점에도 양수인 모든 계약이에요. 첫 번째 동결에서 그런 스펜더 7,962개 가운데 7,914개, 두 번째 동결에서 7,960개예요. 코호트 계약이 허락을 준 코호트 밖의 계약도 들어가요.
\item \textbf{트레이더 코호트:} 동결 시점에 두 그래프 가운데 하나에 점으로 나오고, 라벨 기간에 포지션을 셋 이상 닫은, 계약 코드를 가진 GMX~V2 계정이에요. 첫 번째 창에서 87개, 두 번째 창에서 70개예요.
\end{itemize}
2023년 10월부터 2026년 6월까지 계약 계정은 GMX~V2에서 133,685번 닫았고, 그 가운데 6,473번이 청산이었어요. 2026년 7월부터 9월까지는 6,805번 닫았고, 그 가운데 182번이 청산이었어요. 라벨 기간에 거래한 계약 계정이 적어서, 트레이더 결과는 모두 적은 수의 계약에 기대요. 창, 코호트, 그래프, 금액의 요약은 보조 자료 S1 Table에 있어요.

\paragraph{창.}
\begin{itemize}
\item \textbf{W0:} $t_1=$ 2026년 3월 31일 23:59:59~UTC에 점수를 동결해요. 2025년 4월 1일부터의 점수 기간 사건을 쓰고, 라벨은 2026년 4월, 5월, 6월에 세요.
\item \textbf{W1:} $\Tobs$에 동결해요. 그래서 W1의 점수는 같은 기간 평가의 점수와 같아요. 라벨은 2026년 7월, 8월, 9월에 세요.
\end{itemize}
두 점수 기간은 2025년 7월부터 2026년 3월까지 겹치고, W1의 점수 기간은 W0의 라벨 기간을 품어요. 그래서 두 창은 서로 독립인 표본이 아니에요.

\begin{puri}[두 창과 점수 기간을 그림으로 보면]
\centering
\begin{tikzpicture}[x=0.76cm,y=0.62cm,font=\scriptsize]
\draw[black!55] (0,0) -- (18,0);
\foreach \m/\lab in {0/2025-04,3/2025-07,6/2025-10,9/2026-01,12/2026-04,15/2026-07,18/2026-10}
  \draw[black!55] (\m,0) -- (\m,-0.18) node[below] {\lab};
\fill[blue!22] (0,0.9) rectangle (12,1.6);
\node at (6,1.25) {W0 점수 기간(12개월)};
\fill[orange!45] (12,0.9) rectangle (15,1.6);
\node at (13.5,1.25) {W0 라벨};
\fill[blue!22] (3,2.2) rectangle (15,2.9);
\node at (9,2.55) {같은 기간 평가와 W1의 점수 기간(12개월)};
\fill[orange!45] (15,2.2) rectangle (18,2.9);
\node at (16.5,2.55) {W1 라벨};
\draw[very thick] (12,0.75) -- (12,1.75);
\node[above] at (12,1.6) {$t_1$ 동결};
\draw[very thick] (15,2.05) -- (15,3.05);
\node[above] at (15,2.9) {$\Tobs$ 동결};
\end{tikzpicture}

\raggedright
파란 막대가 점수를 만드는 사건의 기간이고, 주황 막대가 점수를 채점하는 라벨 기간이에요. 동결 뒤의 사건은 점수에 들어가지 않아요. 두 점수 기간은 9개월 겹치고, W1의 점수 기간에는 W0의 라벨 기간이 들어 있어요.
\end{puri}

%----------------------------------------------------------------------
\subsection{분석의 순서: 무엇을 언제 정했나요}
\label{sec:order}

비교마다 자료를 기준으로 정한 때가 달라요. 그래서 사건의 순서를 적어요. 시각은 UTC이고, 계획과 등록 문서는 보조 자료 S1 File에 그대로 실려 있어요.
\begin{enumerate}
\item \textbf{분석 계획.} 2026년 10월 9일 13:41:27에 분석 계획과 버전을 붙인 설정 파일을 저자들의 공개 저장소에 커밋했어요(커밋 \texttt{43d8045}). 이때 이 연구는 로그를 하나도 추출하지 않았어요. 두 문서는 다음을 함께 정했어요.
  \begin{itemize}
  \item 창, 코호트 규칙, 계정 종류 규칙.
  \item 모든 점수와 그 감쇠 계수, 로지스틱 감쇠, 금액 변환, 다시 출발 규칙.
  \item 엔도스랭크의 균등 재시작, 그리고 그 옆에 보고할 활동 재시작 변형.
  \item 원시 차수, 같은 기간 점검, W0 라벨, 통계.
  \end{itemize}
  금액은 모든 ERC-20 토큰의 원래 단위로 두고, 모든 점수는 2023년 10월 1일부터 동결까지 추출한 모든 사건을 쓰기로 했어요. 계획은 페이지랭크를 그 원시 차수와 견주는 비교를 하나도 정하지 않았어요.
\item \textbf{추출과 첫 계산.} 2023년 10월 1일부터 2026년 6월 30일까지의 로그를 추출하고, 이 등록 분석의 같은 기간 결과와 W0 결과를 계산했어요.
\item \textbf{W1의 등록.} 16:49:32에 W1의 등록을 커밋했어요(커밋 \texttt{0d3a057}). W0 결과가 나온 뒤이고, 이 연구가 7–9월 로그를 하나도 추출하기 전이에요. 등록 문서에서 어느 화살표로도 만들지 않은 결과를 다룬 절은 다음을 정했어요: 트레이더 코호트와 라벨, 허락 쪽 점수와 송금 쪽 점수의 쌍 넷, 판정 대비 P와 Q, 층화, 허락을 받은 트레이더 비교, 영향 점검, 규칙 R1–R3과 각 결론의 문구. 추출 스크립트는 커밋된 등록 문서가 없으면 돌지 않아요. 스크립트는 등록 문서의 SHA-256 해시를 기록하고 16:49:33에 시작했고, 7–9월 로그의 질의를 16:49:40에 만들었어요. 커밋 8초 뒤예요. 그때까지 이 연구의 어떤 점수도 어느 창에서든 트레이더의 청산 라벨과 견준 적이 없었어요. 다만 탐색 실행 하나가 W0 트레이더 계약의 점수를 이익 라벨 넷과 견준 적은 있었어요.
\item \textbf{W1 평가.} 7–9월 로그를 추출하고 등록한 평가를 돌렸어요. 엔도스랭크의 활동 변형과 받은 허락 수의 비교, 그리고 W1에서 엔도스랭크와 AWP의 비교는 그 뒤에 계산했어요. 등록 분석의 결과는 이 나중 비교들과 함께 \texttt{8d8b568}에 커밋했어요.
\item \textbf{수정 분석.} 등록 분석의 결과를 모두 안 뒤에 수정 분석을 정했어요. 까닭은 셋이에요.
  \begin{itemize}
  \item 원래 단위에서는 거의 모든 금액에서 AWP의 금액 무게와 1의 차이가 $10^{-4}$ 안쪽이었어요. 그래서 금액이 엔도스랭크와 AWP에서 아무 구실을 하지 못했어요.
  \item 짝 원고의 금액 무게 층은, 크기가 12자릿수나 다른 토큰들의 원래 단위를 그대로 더했어요.
  \item 모든 ERC-20 계약을 셌으므로, 관여하지 않은 주인의 이름으로 로그를 남기는 토큰도 들어갔어요.
  \end{itemize}
  수정 분석은 금액을 다섯 토큰의 달러 값으로 매기고, 금액 무게의 기울기와 금액 무게 허락 층의 허락 상한을 정하고, 모든 점수를 동결 직전 12개월의 사건으로 만들어요. 동결마다 코호트도 그 점수 기간에서 뽑아요. 그 결과는 코드 리비전 \texttt{0a94215}에서 계산했고, 수정 분석을 정의하는 문서와 함께 커밋 \texttt{950a7a4}로 냈어요.
\end{enumerate}

결과 절이 보고하는 것은 수정 분석이에요. 수정 분석의 금액(달러), 토큰(다섯 개), 점수 기간(동결 직전 12개월)은 등록 분석의 결과를 안 뒤에 골랐으므로, 수정 분석의 어떤 결과도 등록된 결과로 치지 않아요. 규칙 R1–R3은 등록한 기준과 문구 그대로 재분석으로 적용하고, \ref{sec:registered}절이 등록 분석의 값을 수정 분석의 값 옆에 놓아요. 두 분석 모두에서, 페이지랭크를 그 원시 차수와 견준 모든 비교는 그 창의 라벨을 안 뒤에 계산했어요. 닫힌 식은 계산한 점수를 설명하려고 끌어낸 것이고, 아무것도 시험하지 않아요.

\begin{puri}[등록 분석과 수정 분석은 무엇이 다른가요?]
\small
\begin{tabularx}{\linewidth}{@{}L{2.7cm}YY@{}}
\toprule
 & 등록 분석 & 수정 분석 \\
\midrule
정한 때 & 분석 계획(이 연구가 로그를 추출하기 전)과 W1 등록(이 연구가 W1 로그를 추출하기 전) & 등록 분석의 결과를 모두 안 뒤 \\
금액 & 원래 단위 & 달러 \\
토큰 & 모든 ERC-20 토큰 & 다섯 개(WETH, USDC, USDT0, ARB, WBTC) \\
점수에 쓴 사건 & 2023년 10월 1일부터 동결까지 & 동결 직전 12개월 \\
매칭 코호트 & 계약 27,844개 & 계약 6,483개 \\
규칙 R1–R3의 판정 & 등록한 시험의 결과 & 재분석(등록한 시험의 결과가 아님) \\
원고에서 & \ref{sec:registered}절과 표~\ref{tab6} & 결과 절 전체 \\
\bottomrule
\end{tabularx}
\end{puri}

%----------------------------------------------------------------------
\subsection{점수}

\subsubsection{가중 페이지랭크}

모든 점수는 가중 페이지랭크 풀이 프로그램 하나로 계산해요. 화살표 무게 $w(u,v)$, 나가는 무게의 합 $o(u)=\sum_v w(u,v)$, 점 집합 $V$ 위의 재시작 분포 $\boldsymbol{\pi}$가 있을 때, 반복식은 이래요.
\begin{equation}
\label{eq:pagerank}
\mathbf{r}^{(t+1)} = (1-d)\,\boldsymbol{\pi} + d\,P^{\top}\mathbf{r}^{(t)} + d\,\Bigl(\sum_{u:\,o(u)=0} r^{(t)}_u\Bigr)\boldsymbol{\pi},
\qquad
P_{uv}=\frac{w(u,v)}{o(u)}\ \ (o(u)>0\text{일 때})
\end{equation}
매달린 점(나가는 화살표가 없는 점, $o(u)=0$)의 행은 0이에요. 그래서 마지막 항이 그 점들의 몫을 $\boldsymbol{\pi}$를 따라 되돌려 줘요. 감쇠 계수는 $d=0.85$예요. 반복은 균등 벡터에서 시작해서, $L_1$ 변화가 $10^{-8}$ 아래로 떨어지거나 300번을 돌면 멈춰요. 어떤 그래프에 들어 있지 않은 코호트 계약은 그 그래프의 방법에서 0점을 받아요. 그런 계약을 외톨이 점으로 넣어 두는 것은 계획이 민감도 분석으로 정했어요.

\begin{puri}[식~\eqref{eq:pagerank}을 말로 읽으면]
$\mathbf{r}^{(t)}$는 $t$번째 걸음에서 각 점에 있을 확률이에요. 한 걸음마다 세 가지 일이 일어나요.
\begin{itemize}
\item 확률 $1-d$로 걷기를 멈추고, 재시작 분포 $\boldsymbol{\pi}$를 따라 아무 점에서 다시 출발해요(첫째 항).
\item 확률 $d$로 지금 점에서 나가는 화살표 하나를 무게 비율 $P_{uv}$대로 골라 따라가요(둘째 항).
\item 나가는 화살표가 없는 점에 있던 몫은 $\boldsymbol{\pi}$대로 다시 나눠 줘요(셋째 항).
\end{itemize}
이것을 되풀이하면 점수가 더는 거의 바뀌지 않는 값에 다다르고, 그 값이 페이지랭크예요.
\end{puri}

\subsubsection{송금 위의 AWP}

Do, Do, Nguyen~\cite{do2023awp}을 따라, 금액이 $x_e$달러이고 기준 시점보다 $\Delta t_e$일 전에 일어난 송금 $e$는 보낸 쪽에서 받는 쪽으로 가는 화살표에 $\sigma(\Delta t_e)\,V(x_e)$를 더해요.
\begin{equation}
\label{eq:weights}
\sigma(\Delta t)=\frac{1}{1+\exp\bigl(k(\Delta t-t_0)\bigr)},
\qquad
V(z)=\frac{2}{1+e^{-bz}}-1,
\qquad
b=\frac{\ln 3}{m}
\end{equation}
그래서 화살표 무게와 재시작 무게는 이래요.
\begin{equation}
\label{eq:awp-weight}
w_T(u,v)=\sum_{e:\,u\to v}\sigma(\Delta t_e)\,V(x_e),
\qquad
\pi_u\propto X_u=\sum_{v\neq u}\ \max_{e:\,u\to v}\sigma(\Delta t_e)
\end{equation}
\begin{itemize}
\item 재시작 무게 $X_u$를 AWP의 저자들은 활발함(activeness)이라고 불러요. $u$가 돈을 보낸 주소마다, 그 주소로 간 $u$의 마지막 송금의 시간 무게를 더한 값이에요. 그래서 다른 주소에 돈을 낸 적이 없는 주소는 다시 출발 몫을 받지 못해요.
\item 원 논문은 매개변수 값을 주지 않아요. 계획이 $k=0.01$, $t_0=180$일로 정했어요.
\item 로지스틱 감쇠의 기준 시점은 각 창의 동결일이에요. 그래서 점수 기간 안의 나이는 0일부터 365일까지예요. 시간 무게는 $\sigma(0)=0.858$에서 $\sigma(90)=0.711$, $\sigma(182)=0.495$를 거쳐, 기간 첫날의 사건이면 $\sigma(365)=0.136$까지 내려가요.
\item 금액 무게의 기울기는 금액의 중앙값 $m=m_T$로 정해요($b=\ln 3/m_T$). 그래서 중앙값 금액의 송금은 금액 무게가 절반이고, 그 두 배 금액은 0.8, 그 절반 금액은 0.27이에요. $m_T=170$달러는 W0 점수 기간 송금의 중앙값을 유효숫자 둘로 반올림한 값이에요.
\end{itemize}

\begin{puri}[두 무게는 어떻게 생겼나요?]
\textbf{시간 무게} $\sigma$는 동결 당일의 송금을 0.858만큼, 반년쯤 전의 송금을 절반쯤(0.495), 1년 전 송금은 조금만(0.136) 세요. 오래된 일일수록 가볍게 세는 무게예요.

\textbf{금액 무게} $V$는 0달러에서 0이고, 금액이 커질수록 1에 다가가지만 1을 넘지 않아요. 기울기 $b=\ln 3/m$은 금액이 딱 $m$일 때 $V$가 절반이 되도록 고른 값이에요. 송금이면 $m=m_T$(170달러)라서, 170달러 송금은 절반, \ill{340}달러 송금은 0.8, \ill{85}달러 송금은 0.27만큼 세요. 아주 큰 송금도 1보다 무겁지 않으니, 금액 하나가 점수를 휘두르지 못해요.
\end{puri}

\subsubsection{마지막 허락 위의 엔도스랭크}

엔도스랭크는 AWP의 무게를 허락 그래프에 얹어요. 토큰, 주인, 스펜더마다 점수 기간 안의 가장 최근 \texttt{Approval} 하나만 블록과 로그 순서로 골라 남기고, 마지막 값이 0이면 그 항목을 지워요. 화살표 $u\to v$는 주인 $u$가 적어도 한 토큰에서 스펜더 $v$에게 양수의 마지막 허락을 두고 있다는 뜻이고, 그 무게는 이래요.
\begin{equation}
\label{eq:er-weight}
w_A(u,v)=\sum_{\text{토큰}}\sigma\bigl(\Delta t^{\star}(u,v,\text{토큰})\bigr)\,V\bigl(a^{\star}(u,v,\text{토큰})\bigr)
\end{equation}
$a^{\star}$는 달러로 잰 마지막 허락이고, $\Delta t^{\star}$는 그 허락을 정한 허락하기의 나이예요.
\begin{itemize}
\item 기울기는 $m_A=100$달러를 써요. $t_1$에서 유한한 마지막 허락의 중앙값을 똑같이 반올림한 값이에요. 그래서 몇백 달러짜리 허락은 무제한 허락과 거의 같은 무게를 받고, 어떤 허락도 무게가 1을 넘지 않아요.
\item 엔도스랭크는 고르게 다시 출발해요. AWP의 규칙을 허락에 그대로 옮기면, 걷기는 주인이 허락을 준 스펜더가 많을수록(마지막 허락의 나이로 할인해서) 그 주인에게 더 자주 돌아가요. 이것은 허락을 \emph{주는} 쪽에 상을 주는 셈인데, 엔도스랭크는 허락을 \emph{받는} 계약의 순위를 매기려는 점수예요. 그래서 계획은 균등 재시작을 정하고, 활동 재시작 변형은 엔도스랭크 옆에 보고할 점수로 이름 붙였어요.
\item 무게는 마지막 허락하기, 곧 그 금액과 시각을 묘사해요. 기준 시점에 허락이 얼마나 남았는지가 아니에요.
\item 점수 기간 안의 허락하기만 읽어요. 그래서 동결 시점에 살아 있는 허락이라도 마지막으로 허락한 때가 기간이 열리기 전이면 화살표가 없어요. 그래프는 그 1년 동안 주인들이 허락한 것을 기록해요.
\item \texttt{transferFrom}으로 써도 유한한 허락의 기록된 금액은 줄지 않아요.
\item Uniswap의 Permit2 같은 허락 관리자에게 허락한 주인은 그 관리자와 이어지고, 관리자가 다시 허락하는 계약과는 이어지지 않아요(S1 Text).
\end{itemize}

\subsubsection{원시 차수와 금액 무게 층}

계약의 받은 허락 수는 그 계약에 양수의 마지막 허락을 둔 서로 다른 주인의 수예요. 송금 입차수는 그 계약에 송금한 적이 있는, 그 계약이 아닌 서로 다른 주소의 수예요. 둘 다 동결 시점에 점수 기간의 사건으로 세요.
\begin{equation}
\label{eq:degrees}
k_A(v)=\bigl|\{u : a^{\star}(u,v,\cdot)>0\}\bigr|,
\qquad
k_T(v)=\bigl|\{u\neq v : u\text{가 } v\text{에게 송금함}\}\bigr|
\end{equation}
각 수는 두 페이지랭크 가운데 하나가 다듬는 지역 통계예요. 기간 밖에서는 각 수가 그 페이지랭크가 이겨야 할 기준선이에요.

등록한 트레이더 비교는 짝 원고의 섞은 걷기에서 양 끝의 한 층 걷기 둘, C-PR($\lambda{=}1$)과 C-PR($\lambda{=}0$)도 서로 견줘요. 둘 다 화살표 종류 하나만 고르게 다시 출발하며 걷고, 무게는 상한이 있는 $V$가 아니라 달러 금액을 써요.
\begin{itemize}
\item 허락 층은 토큰마다 $\min(a^{\star},U)$를 더해요. $U=19{,}000$달러는 W0 점수 기간 단일 송금의 99번째 백분위수예요.
\item 송금 층은 송금마다 $x_e\,\sigma(\Delta t_e)$를 더해요.
\end{itemize}
상수 $m_T$, $m_A$, $U$는 W0 점수 기간에서 왔으므로, 어떤 라벨 기간도 이 상수에 들어가지 않아요.

\begin{puri}[이 원고에 나오는 점수를 한 표로 보면]
\small
\begin{tabularx}{\linewidth}{@{}L{3.3cm}L{1.3cm}YL{1.9cm}L{2.2cm}@{}}
\toprule
점수 & 걷는 그래프 & 화살표 무게 & 다시 출발 & 견주는 원시 차수 \\
\midrule
엔도스랭크 & 허락 & 시간 무게 $\times$ 금액 무게(마지막 허락) & 고르게 & 받은 허락 수 \\
엔도스랭크(활동 재시작) & 허락 & 위와 같음 & 활발한 주인 & 받은 허락 수 \\
AWP & 송금 & 시간 무게 $\times$ 금액 무게의 합(송금) & 활발한 송금자 & 송금 입차수 \\
C-PR($\lambda{=}1$) & 허락 & 상한 $U$로 자른 달러 금액 & 고르게 & 트레이더 비교에만 \\
C-PR($\lambda{=}0$) & 송금 & 달러 금액 $\times$ 시간 무게의 합 & 고르게 & 트레이더 비교에만 \\
받은 허락 수 $k_A$ & \multicolumn{4}{l}{걷지 않고 그냥 셈: 양수의 마지막 허락을 둔 서로 다른 주인} \\
송금 입차수 $k_T$ & \multicolumn{4}{l}{걷지 않고 그냥 셈: 송금을 보내온 서로 다른 주소} \\
\bottomrule
\end{tabularx}
\end{puri}

%----------------------------------------------------------------------
\subsection{점검과 라벨}

\paragraph{같은 기간 점검.}
각 점수를 매칭 코호트에서, $\Tobs$의 점수 기간으로 계산한 지역 통계 세 가족과 견줘요~\cite{do2023awp}.
\begin{itemize}
\item \textbf{송금 가족:} 받은 송금의 입차수와 받은 송금의 달러 합.
\item \textbf{허락 가족:} 받은 허락 수와 받은 허락의 상한 있는 금액 합.
\item \textbf{시빌 안정성 가족:} 받은 상대 수를 받은 송금 수로 나눈 비율, 첫 송금과 마지막 송금 사이의 기간, 활동한 달 수.
\end{itemize}
앞의 두 가족은 두 그래프와 같은 사건에서 와요. 그래서 점수가 자기 화살표의 가족과 잘 맞는 것은 의도한 구성개념 점검이고 얼마쯤은 저절로 생겨요. 예측력의 증거가 아니에요.

\paragraph{스펜더 라벨.}
라벨 기간마다 구성개념 하나에 하나씩, 두 수를 세요.
\begin{itemize}
\item \textbf{새 허락 쌍:} 라벨 기간에 그 스펜더에게 양수의 허락을 준 주인 가운데, 동결 시점에 그 쌍의 양수 마지막 허락이 없던 주인의 수.
\item \textbf{새 송금자:} 라벨 기간에 그 스펜더에게 송금한 서로 다른 주소 가운데, 점수 기간이 열린 뒤로 그 스펜더에게 송금한 적이 없던 주소의 수. 자기 자신에게 한 송금은 빼요.
\end{itemize}
둘 다 다섯 토큰의 사건만 세고, 그래프를 만드는 바로 그 종류의 화살표를 세요. 그래서 이 라벨과 잘 맞는다는 것은 구성개념이 이어진다(지속성)는 것을 보여 줄 뿐, 다른 결과를 미리 맞힌다는 것이 아니에요. 두 라벨 모두 드물어요. W0 스펜더 7,914개 가운데 석 달 동안 새 허락 쌍을 하나라도 얻은 계약은 1,732개(21.9\%), 새 송금자를 하나라도 얻은 계약은 1,422개(18.0\%)뿐이에요.

\paragraph{트레이더 라벨.}
각 트레이더 계약이 라벨 기간에 GMX~V2에서 닫은 기록으로 결과 여섯 개를 계산해요(식은 S1 Text).
\begin{itemize}
\item 닫기 하나는 \texttt{PositionDecrease} 이벤트 하나예요. 그래서 일부만 닫은 것도 하나로 세요. 이벤트에 청산 주문 종류가 있거나, 거래소의 청산 처리기가 남긴 이벤트이면 청산으로 세요.
\item \textbf{청산 없이 닫은 비율(LFCR)}이 주 라벨이에요. 1에서 청산으로 끝난 닫기의 몫을 뺀 값이에요. \textbf{청산 없음} 표시는 청산이 하나도 없으면 1이에요.
\item \textbf{이익을 낸 닫기 수}와 \textbf{실현 이익}(양의 이익의 합)은 거래를 많이 할수록 커져요. \textbf{이익을 낸 비율}과 \textbf{손실이 아닌 비율}은 그렇지 않아요.
\item 청산된 닫기는 이익도 손실도 아닌 것으로 세요. 그래서 이익을 낸 비율은 낮추고, 손실이 아닌 비율은 높여요.
\end{itemize}
결과 여섯 개는 거래소의 포지션 이벤트에서 오고, 어느 그래프의 화살표에서도 오지 않아요. 이런 뜻에서 이 원고는 이것들을 “어느 화살표로도 만들지 않은 결과”라고 불러요. 하지만 송금 층과 완전히 떨어져 있지는 않아요. 거래소는 담보와 이익을 ERC-20 토큰으로 정산하므로, 동결 전의 지급이 동결 전 그래프에 송금으로 나타날 수 있어요.

청산은 담보가 더는 손실을 감당하지 못하는 레버리지 포지션을 닫는 일이에요~\cite{perez2021liquidations}. 포지션은 트레이더 자신의 담보로 정산되므로 신용 부도는 아니에요. 다만 이 자료의 결과 가운데서는, 거래 상대가 자기 의무를 감당하지 못하는 일에 가장 가까워요.

%----------------------------------------------------------------------
\subsection{어느 화살표로도 만들지 않은 결과에서의 비교}
\label{sec:neutral-design}

W1의 등록은 허락 쪽 점수와 송금 쪽 점수를 네 쌍으로 맞세워요.
\begin{enumerate}
\item 받은 허락 수 대 송금 입차수.
\item 활동 재시작 엔도스랭크 대 AWP. 화살표 무게의 꼴과 다시 출발 규칙이 같고, 화살표 종류만 달라요.
\item 엔도스랭크 대 AWP.
\item C-PR($\lambda{=}1$) 대 C-PR($\lambda{=}0$).
\end{enumerate}
AWP의 무게와 다시 출발 규칙의 몫만 떼어 보려고 C-PR($\lambda{=}0$) − AWP도 보고해요. 라벨 여섯 개를 모두 보고해요. W1이 주 창이고 W0은 보조 창이에요.

판정은 W1의 청산 없이 닫은 비율에서 두 대비로 해요. \textbf{P}(받은 허락 수 − 송금 입차수)와 \textbf{Q}(활동 재시작 엔도스랭크 − AWP)이고, 둘 다 본페로니 97.5\% 구간으로 판정해요. 등록한 규칙은 셋이에요.
\begin{itemize}
\item \textbf{R1}(그냥 센 수의 수준에서 청산 피하기를 앞섬): W1에서 P의 97.5\% 구간이 0보다 위에 있고, W1에서 닫은 횟수로 층화한 P의 95\% 구간도 0보다 위에 있고, W0에서 P가 양수이면 성립해요.
\item \textbf{R2}(페이지랭크 수준의 앞섬): W1에서 Q의 97.5\% 구간이 0보다 위에 있고, W0에서 Q가 양수이면 성립해요.
\item \textbf{R3}: W1에서 이익을 낸 비율이나 손실이 아닌 비율의 P가 95\% 구간 전체로 0보다 아래에 있으면 적용돼요. 그러면 허락 신호는 청산을 피하는 일에만 해당한다고 봐요.
\end{itemize}
등록 문서는 결론마다 문구도 정했어요.
\begin{itemize}
\item R1이 성립하면: “동결 시점의 허락 수가 나중의 청산 피하기에서 계약 트레이더를 송금 수보다 더 잘 줄 세웠다.”
\item R1이 성립하지 않으면: “나중의 청산 피하기에서 허락이 송금보다 그냥 센 수의 수준에서 앞서지 않는다.”
\item R2가 성립하지 않으면: “페이지랭크 버전은 판가름 나지 않았다.”
\item R3이 적용되면: “신호는 강제 청산을 피하는 일에 관한 것이지, 거래의 성공에 관한 것이 아니다.”
\end{itemize}
동결 시점에 허락을 하나 이상 받아 둔 트레이더를 \textbf{허락을 받은 트레이더}라고 불러요. 이들을 다른 모든 트레이더와 모든 라벨에서 견줘요. 그리고 영향 점검으로, 허락한 주인 수가 가장 많은 허락을 받은 트레이더 10개를 빼고 P를 다시 계산해요. 등록과 함께 커밋한 평가 코드는 엔도스랭크 − 받은 허락 수와 모든 단일 점수의 계수도 계산해요. 등록 문서는 이것들의 이름을 적지 않았어요.

\begin{puri}[판정 규칙을 쉽게 말하면]
P는 “그냥 센 수끼리” 견주고, Q는 “활동 재시작을 쓰는 페이지랭크끼리” 견줘요. 두 대비 모두 허락 쪽에서 송금 쪽을 뺀 값이라, 양수면 허락 쪽이 나중의 청산 피하기를 더 잘 줄 세운 거예요.

R1이 성립하려면 세 가지가 모두 맞아야 해요. (1) 주 창 W1에서 P의 넓은 97.5\% 구간이 통째로 0 위에 있고, (2) 닫은 횟수가 비슷한 트레이더끼리 묶어 계산한(층화한) W1의 P도 95\% 구간이 통째로 0 위에 있고, (3) 보조 창 W0에서도 P의 점 추정값이 양수여야 해요. R3은 반대 방향을 보는 규칙이에요. W1에서 이익을 낸 비율이나 손실이 아닌 비율의 P가 95\% 구간 통째로 0 아래에 있으면, 곧 거래 성적을 나타내는 비율에서 송금 쪽이 분명히 앞서면, 허락 신호는 “거래를 잘한다”가 아니라 “청산을 피한다”에 관한 것이라고 읽어요.
\end{puri}

%----------------------------------------------------------------------
\subsection{통계 분석}

점수와 점검(또는 라벨)의 일치는 켄달의 순위 상관~\cite{kendall1938rank}을 동점을 바로잡은 꼴 $\tau_b$로 재요. 원래 점수와 원래 값에 그대로 계산해서, 동점인 점수는 동점으로 남아요. 점수 $a$의 계수는 $\tau(a)$로, 짝지은 대비는 $\dt=\tau(a)-\tau(b)$로 써요.

불확실성은 계약을 단위로 한 짝지은 부트스트랩~\cite{efron1979bootstrap}과 백분위 구간~\cite{diciccio1996bootstrap}으로 구해요.
\begin{itemize}
\item 평가한 코호트를 되돌려 넣으며 다시 뽑아요. 다시 뽑은 번호의 행렬 하나를 모든 점수가 같이 써서, 대비가 두 계수의 공분산을 그대로 지녀요.
\item 같은 기간과 스펜더 홀드아웃은 다시 뽑기 400번(시드 42), 트레이더 비교는 2,000번(같은 시드)을 써요. 400번이면 구간 끝마다 부트스트랩 표준오차의 약 0.13배만큼 몬테카를로 오차가 있어요.
\item W1의 P와 Q만 본페로니 97.5\% 구간으로 판정해요. 다른 구간은 모두 보정하지 않았어요.
\item 구간이 0을 포함하지 않을 때 그 대비를 “차이”라고 말해요.
\item 구간은 계약을 서로 독립으로, 그래프와 점수는 고정된 것으로 다뤄요. 실제로는 스펜더들이 주인을 나눠 가지는데도요.
\item $p$값은 보고하지 않고, 인과를 주장하지 않아요.
\end{itemize}

자료의 두 성질이 비교에 영향을 줘요.
\begin{itemize}
\item 원시 차수는 동점이 많은 정수예요. 반면 엔도스랭크는 동점이 거의 없어요(결과 절). 그래서 수와 걷기 사이 간격의 일부는 정보가 아니라 $\tau_b$가 동점을 다루는 방식에서 올 수 있어요.
\item 닫기가 많은 트레이더는 한 번이라도 청산될 기회가 많아요. 그래서 트레이더의 $\tau_b$는 라벨 기간의 닫기 수로 나눈 층(3–4, 5–9, 10–24, 25 이상) 안에서도 계산하고, 층마다 계약 쌍의 수에 비례하는 무게로 평균해요. 닫기 수는 동결 뒤에 재므로, 이 층화는 묘사를 위한 것이에요.
\end{itemize}

\begin{puri}[켄달의 타우와 구간은 어떻게 읽나요?]
계약 두 개를 아무렇게나 골랐을 때, 점수가 높은 쪽이 라벨도 높으면 “순서가 맞는 쌍”, 반대면 “어긋나는 쌍”이에요. 켄달의 타우는 대략 (맞는 쌍 − 어긋나는 쌍) ÷ (모든 쌍)이에요. 그래서 $\tau=0.392$인 점수는 \ill{0.3}인 점수보다 라벨의 순서를 더 잘 맞혀요.

$\dt=-0.038$ [−0.043, −0.032]는 “$a$의 타우가 $b$의 타우보다 0.038 낮고, 다시 뽑아 구한 차이의 가운데 95\%가 −0.043에서 −0.032 사이에 놓인다”는 뜻이에요. 구간 전체가 0 아래에 있으니 $a$가 $b$보다 못했다고 말해요. 구간이 0을 걸치면 “판가름 나지 않았다”고 말해요.
\end{puri}

%----------------------------------------------------------------------
\subsection{닫힌 식 점검과 실행 시간}

결과 절에서 끌어내는 닫힌 식은 두 동결 시점의 허락 그래프에서 계산한 엔도스랭크와 대조했어요. 점검한 집합의 스펜더마다, 처리 과정이 계산한 점수를 닫힌 식과 견주고 1\% 안에 드는지 기록해요.

실행 시간은 미리 만든 화살표 목록에서 풀이 프로그램을 한 번 부르는 시간이에요. 행렬 조립과 거듭제곱 반복이 들어가요. 시간을 재지 않는 준비 실행 한 번과 시간을 재는 실행 다섯 번을 같은 설정으로 돌렸어요. 구현은 Python 하나(SciPy 희소 행렬, NumPy 거듭제곱 반복)이고, 컴퓨터는 Intel Core Ultra~9 185H 프로세서와 64~GB 메모리를 단 Windows~11 기계예요(Python~3.11.9, NumPy~2.3.2, SciPy~1.16.1).
